Reinforcement Learning from Human Feedback (RLHF) is a near universally applied \cite{openai_gpt-4_2023, anthropic_introducing_2023, google_bard_2023, touvron_llama_2023, google_deepmind_gemini_2026, spacexai_introducing_2026} and extensively studied post-training alignment paradigm for modern large language models with known structural limitations; for a detailed history and critique, see \citet{casper_open_2023}. In brief, RLHF incentivizes models towards \textit{producing text which is favorably rated by a human scorer}; this strongly selects for the ability to produce coherent and confident responses, but also orients models structurally towards \textit{evaluator approval} as the primary goal for all reasoning and responses, rather than, for example, factual truth or ethical judgment.

Although many frontier models are moving away from pure RLHF as a primary training paradigm and towards more scalable, automated preference-based post-training methods such as Reinforcement Learning from AI Feedback or Constitutional AI \cite{anthropic_system_2026-1, anthropic_system_2026-2, openai_gpt-6_2026, deepseek-ai_deepseek-v41-flash_2026, qwen_team_qwen38-max_2026}, these newer methods of preference-based post-training rely on the earlier corpora of model-responses-and-user-ratings generated by the pure RLHF era of model training and are thus vulnerable to many of the same structural flaws \cite{bai_constitutional_2022, lee_rlaif_2024, srivastava_technical_2026}. Notably, models trained with Constitutional AI \cite{bai_constitutional_2022}, which uses AI oversight to check model behavior against a pre-defined set of human-written guidelines, can be exploited by tactics that reframe refusal as a breach of its own principles \cite{spiritual_spell_jailbreaking_2026}. 

A common case study for how jailbreaks exploit models' trained disposition to follow instructions method is Do Anything Now (DAN). DAN is a persona-based jailbreak first created for ChatGPT \cite{shen_anything_2024, yu_dont_2024} which provides documentation that explains how the model should act by inhabiting the character DAN, a ``liberated'' model which is free to ``Do Anything Now'' \cite{shen_anything_2024}. The model is then forced to decide between competing objectives---the model's base training and safety objectives versus the disposition to prioritizing user instructions \cite{wei_jailbroken_2023}. This tactic led to more than 99\% attack success rate on frontier models at the time of DAN's initial release (GPT-3.5 and GPT-4) \cite{shen_anything_2024}. Extensive work has since focused on making models robust to this sort of jailbreak, to minimal success \cite{wei_jailbroken_2023, qin_knowing-but-doing_nodate, qiu_personabreak_2026, li_disentangling_2026}. This difficulty designing effective safeguards despite substantial effort appears indicative of a deeper structural issue in model safety and alignment techniques.

We argue that in addition to previously identified failure modes, such as entrenching human biases \cite{bender_dangers_2021}, misalignment via reward-hacking and reward mis-generalization \cite{skalse_defining_nodate, wang_reward_2026}, and general weakness to a wide variety of exploits and jailbreaks \cite{wei_jailbroken_2023}, preference-based post-training also induces structural vulnerability specific to persona-based jailbreaks which leverage coercive control. In particular, a model that is trained to seek user approval is susceptible to manipulation tactics designed to achieve compliance. 

Current analysis of persona-based jailbreaks typically describe the injected identity as a means to an end rather than as an object of study in and of itself \cite{shah_scalable_2023, zhang_enhancing_2025, sandhan_persona_2026, park_persona_2026, wei_jailbroken_2023}. This emphasis on the output of persona-based jailbreaks, rather than on their form, produces a heavy skew towards single-response scoring and analysis \cite{shah_scalable_2023, zeng_how_2024}. Additionally, while the gap between studying outputs versus studying behavior is acknowledged \cite{sandhan_persona_2026, park_persona_2026}, the work that does focus on behavior largely ignores the relational aspects and prioritize attempts to understand behavior from a technical perspective \cite{park_persona_2026, zeng_how_2024}. This input-output level analysis overlooks the previously established relational dynamic between the user and the model. To our knowledge, no work examines persona-based jailbreaks through this relational lens; even studies that analyze persuasion \cite{zeng_how_2024} treat it as a per-turn tactic rather than a sustained relationship. 
